\subsection{$m^2$ of same-event and mixed-event matches per centrality}
\label{sec:m2}

This section uses the 1M-event sample of Sec.~\ref{sec:me1m} (\fpath{StRoot/StTofMEMaker} version 3,
\fpath{output/me_v3.root}). It answers the question: what does the $m^2$ distribution of
random TPC--BTOF matches look like, and how much of the measured $m^2$ distribution does it account for, in each
momentum bin and centrality class? Numbers: \fpath{logs/m2_numbers_v3.txt} from \fpath{scripts/make_m2.C}.

\subsubsection{Method}
\paragraph{Centrality.} Centrality is taken from StRefMultCorr (Au+Au 19.6\,GeV 2019 parameters: $V_z$- and
luminosity-corrected refMult, nine classes from 0--5\% to 70--80\%). Of the \mtEvtSelVthree{} selected events,
\mtEvtPileVthree{} are flagged as pile-up by its refMult--nBTOFMatch cut, \mtEvtPeriphVthree{} are beyond 80\%,
and \mtEvtNoTimeVthree{} have no BTOF hit with a time (no valid start time). The remaining \mtEvtUsedVthree{} events
enter the $m^2$ histograms. Bad runs are excluded at the run selection. The mixing classes are unchanged (raw refMult).

\paragraph{Hit times.} A picoDst BTOF hit carries only its cell. A time of flight exists for a hit only if a picoDst
track was matched to it in production (\fcmd{StPicoBTofPidTraits::btof()}, with $\beta>0$). On average
\mtTimedFracVthree\% of the hits per event have a time. Each cell therefore has one of three states: no hit, hit with
time $t$, hit without time.

\paragraph{$\beta$ and $m^2$.} For each projected QA track the hit is chosen as in production: the projected cell,
otherwise the nearer, then the farther neighbour, preferring hits that have a time (Sec.~\ref{sec:me}). The path length is
$L=$ \fcmd{tofPathLength}(primary vertex, projected point, helix curvature), the same function as
\fcmd{StBTofMatchMaker}, and
\begin{equation}
\beta = \frac{L}{c\,t},\qquad m^2 = p^2\left(\frac{1}{\beta^2}-1\right),
\end{equation}
with $p$ the primary momentum. For tracks whose own production hit is chosen, $L$ agrees with production
($L-\beta_{\rm prod}c\,t_{\rm prod}$: mean \mtDLMeanVthree\,cm, RMS \mtDLRmsVthree\,cm) and $m^2$ differs from the
production value by \mtDmRmsVthree\,GeV$^2$/$c^4$ RMS. The same-event (SE) $m^2$ is therefore the production $m^2$
of the projected tracks.

\paragraph{Mixed event with overlay.} A random hit only produces a wrong $m^2$ when it is chosen \emph{instead of}
the track's own hit. The mixed event (ME) is therefore built as an overlay: the hits of a partner event (same run,
$V_z$ and refMult class, up to 5 partners) are added on top of the hits of the track's own event, and the hit choice
above is redone. When a partner hit is chosen (projected cell free in this event, or same cell with an earlier
partner time), its time is combined with the track's $L$ and $p$ to give an $m^2$. Without the overlay, every
partner hit present in the matching cells would count. That version (kept for reference, \fcmd{hM2ME}) exceeds the
SE yield in the $\pi$--K and K--p valleys by a factor \mtValleyPlainMinVthree--\mtValleyPlainMaxVthree{} in 0--5\%
($0.4<p<1.5$\,\GeVc), which is impossible. The overlay keeps \mtShadowFracVthree\% of those entries.

\paragraph{Normalization.} No fit is used. In each centrality class $c$ and fine momentum bin $j$ (0.05\,\GeVc)
\begin{equation}
\frac{dN_{\rm rand}}{dm^2}(c,j) = \frac{dN_{\rm ME}}{dm^2}(c,j)\;\frac{N_{\rm proj}(c,j)}{N_{\rm trials}(c,j)}\;
\frac{n_{\rm t}+n_{\rm u}}{n_{\rm t}},
\end{equation}
where $N_{\rm proj}$ counts projected tracks, $N_{\rm trials}$ the (track, partner) pairs, and $n_{\rm t}$ ($n_{\rm u}$)
the trials whose chosen partner hit has (has no) time. The last factor, \mtTimedScaleAllVthree{} on average
(\mtTimedScaleCaVthree{} in 70--80\%, \mtTimedScaleCiVthree{} in 0--5\%), restores the random hits that have no stored
time: only \mtMETimedFracVthree\% of the partner hits chosen in the mixed events have a time.
Coarse momentum bins are sums of the normalized fine bins.

\choice{All BTOF hits are included in the random-match estimate (user's decision). Random hits without a stored time
are assumed to have the same time distribution as random hits with a time.}
\caveat{Hits without a time are those no picoDst track was matched to: secondaries, particles not reconstructed in the
TPC, noise. Their times are probably later on average than those of matched hits, so the shape of the random $m^2$
distribution may be biased towards small $m^2$; its normalization would be lower by a factor
\mtTimedScaleAllVthree{} if only timed hits were counted.}

\subsubsection{Results}
\caveat{The overlay used here always keeps the track's own hit in the projected cell. Production instead picks the
hit with the largest ToT among the projected and neighbouring cells, so the random matches below are underestimated.
Sec.~\ref{sec:bkg} studies this and the late background the picoDst cannot describe.}
Figures~\ref{fig:m2_v3_cent8}--\ref{fig:m2_v3_allcent} overlay the normalized ME distribution on the SE $m^2$
distribution in nine momentum bins for 0--5\%, 30--40\%, 70--80\% and 0--80\%; the other centralities are in
Appendix~\ref{app:m2}. Table~\ref{tab:m2contam} and Fig.~\ref{fig:m2_v3_frand} give the random fraction under the
$\pi$, K and p peaks.

\begin{itemize}
\item \textbf{Size.} Integrated over $m^2$ and $0.2<p<3$\,\GeVc, random matches are \mtMENormFracVthree\% of the SE
  matches in 0--80\%, \mtMENormFracCiVthree\% in 0--5\%, \mtMENormFracCeVthree\% in 30--40\% and \mtMENormFracCaVthree\% in 70--80\%.
  This is smaller than $f_{\rm rand}$ of Sec.~\ref{sec:me1m}, because the track's own hit is usually chosen first.
\item \textbf{Under the peaks (0--5\%).} At $0.2<p<0.4$\,\GeVc{} the random fraction is \mtFrPiLoCiVthree\% ($\pi$),
  \mtFrKLoCiVthree\% (K) and \mtFrPLoCiVthree\% (p); at $0.8<p<1.0$\,\GeVc{} it is \mtFrPiMidCiVthree\%, \mtFrKMidCiVthree\% and
  \mtFrPMidCiVthree\%. In 30--40\% the values are about three times smaller (\mtFrPiMidCeVthree\%, \mtFrKMidCeVthree\%,
  \mtFrPMidCeVthree\% at $0.8<p<1.0$\,\GeVc), and in 70--80\% they are below 0.4\%. The kaon fraction stays roughly flat
  with $p$, because random $m^2$ values pile up between the pion and proton peaks.
\item \textbf{Shape (read by eye).} The random distribution peaks just above $m^2=0$ and falls steadily towards large
  $m^2$. At low momentum it follows the SE distribution in the $\pi$--K valley; at higher momentum it lies well below
  the SE distribution everywhere except in the valleys.
\item \textbf{Consistency.} In 0--5\%, $0.4<p<1.5$\,\GeVc, the largest ME/SE ratio in any $m^2$ bin (0.04\,GeV$^2$/$c^4$,
  at least 50 SE counts) is \mtValleyAllVthree. In the lowest momentum bin (0.2--0.4\,\GeVc) it is 1.21
  (Table~\ref{tab:m2contam}), so the estimate is somewhat too large there. The excess probably comes from the all-hits
  scaling, or from the less precise projection at low momentum (Sec.~\ref{sec:me}).
\item \textbf{Tail above the proton.} For $1.6<m^2<2.8$\,GeV$^2$/$c^4$ (between p and d) the SE yield is \mtSbRatioAllVthree{}
  times the random estimate in 0--80\% (2--9 times in 0--5\% above 0.6\,\GeVc, Table~\ref{tab:m2contam}). Random matches account for only a small
  part of this tail.
\end{itemize}

\caveat{Most of the SE yield between the proton and deuteron peaks is not from random hits. Candidates: decays in
flight (the track's $p$ does not belong to the particle that fired the cell), wrong primary vertex or out-of-time
pile-up, and wrongly matched neighbour hits of the same event. Not studied.}

\begin{table}[h]
\centering
\caption{Random-match fraction (\%) of the SE yield in the $\pi$ ($-0.1<m^2<0.1$), K ($0.17<m^2<0.33$) and p
($0.6<m^2<1.2$\,GeV$^2$/$c^4$) windows, SE/ME ratio between the p and d peaks, and the largest ME/SE ratio over $m^2$ bins
with $\ge50$ SE counts in $-0.3<m^2<3$. ME: overlay, all hits. Statistical errors not shown; they are large in
70--80\% above 1\,\GeVc. Generated: \fpath{tables/m2_contam_v3.tex}.}
\label{tab:m2contam}
\small
\input{tables/m2_contam_v3.tex}
\end{table}

\begin{figure}[H]
\centering
\figpanel{1.0}{m2_v3_frand}{Random-match fraction under the $\pi$ (left), K (middle) and p (right) peaks vs momentum
for 0--5\% (red), 10--20\% (orange), 30--40\% (green), 50--60\% (blue) and 70--80\% (violet). Points with fewer than
100 SE counts in the window are omitted. ME: overlay, all hits.}
\caption{Random fraction under the PID peaks.}
\end{figure}

\newcommand{\mtwocap}[1]{$m^2$ of projected QA tracks, #1, nine momentum bins (log scale, 0.04\,GeV$^2$/$c^4$ bins).
Black points: same event (cell$\pm$1 rule, hit with time). Red filled: mixed event (overlay, all hits) normalized per
projected track. Statistical errors on the SE points only.}

\clearpage
\begin{figure}[H]
\centering
\figpanel{0.82}{m2_v3_cent8}{\mtwocap{0--5\%}}
\caption{SE vs ME $m^2$, 0--5\%.}
\end{figure}

\begin{figure}[H]
\centering
\figpanel{0.82}{m2_v3_cent4}{\mtwocap{30--40\%}}
\caption{SE vs ME $m^2$, 30--40\%.}
\end{figure}

\begin{figure}[H]
\centering
\figpanel{0.82}{m2_v3_cent0}{\mtwocap{70--80\%}}
\caption{SE vs ME $m^2$, 70--80\%.}
\end{figure}

\begin{figure}[H]
\centering
\figpanel{0.82}{m2_v3_allcent}{\mtwocap{0--80\% (sum of the nine classes)}}
\caption{SE vs ME $m^2$, 0--80\%.}
\end{figure}
